% ═════════════════════════════════════════════════════════════════════════════
%  chapitres/introduction.tex
% ═════════════════════════════════════════════════════════════════════════════
\chapter*{Introduction Générale}
\addcontentsline{toc}{chapter}{Introduction Générale}
\markboth{Introduction Générale}{Introduction Générale}

% ─────────────────────────────────────────────────────────────────────────────
\section*{Contexte et Problématique}
% ─────────────────────────────────────────────────────────────────────────────

L'industrie du logiciel connaît depuis une décennie une transformation profonde,
portée par deux évolutions majeures. D'une part, le \textbf{passage aux architectures
microservices} a fragmenté les applications monolithiques en un ensemble de services
indépendants, chacun pouvant être développé, déployé et mis à l'échelle
individuellement. D'autre part, l'adoption généralisée des pratiques
\textbf{DevOps} a raccourci drastiquement les cycles de livraison, permettant de
passer de plusieurs déploiements par an à plusieurs déploiements par jour.

Si ces deux tendances ont considérablement amélioré l'agilité et la résilience des
systèmes, elles ont simultanément introduit une \textbf{surface d'attaque étendue}
et difficile à maîtriser. Un système composé de dizaines, voire de centaines de
microservices communicants génère autant de canaux de communication potentiellement
vulnérables, et la vitesse de déploiement laisse peu de place aux audits de sécurité
traditionnels, souvent relégués en fin de cycle.


  \textbf{Constat fondamental :} Selon le rapport \textit{State of DevOps} 2023,
  plus de 60\,\% des incidents de sécurité en production dans les environnements
  cloud-natifs sont liés à des communications interservices non sécurisées ou à
  des configurations d'accès insuffisamment contrôlées.


Face à ce constat, la communauté technique a progressivement convergé vers deux
approches complémentaires. Le paradigme \textbf{DevSecOps} préconise d'intégrer
la sécurité dès les premières phases du développement ; le principe du
\textit{Shift Left Security} en automatisant les contrôles de sécurité au sein
des pipelines CI/CD. Le \textbf{Service Mesh}, quant à lui, propose une couche
d'infrastructure dédiée à la gestion des communications interservices, offrant
nativement des mécanismes de chiffrement, d'authentification et d'observabilité
sans modifier le code applicatif.

La problématique centrale de ce projet est donc la suivante :


  \textit{Dans quelle mesure l'intégration d'un Service Mesh au sein d'une
  démarche DevSecOps permet-elle de renforcer significativement la sécurité
  des communications interservices dans les architectures microservices, et
  quelles sont les conditions d'une telle intégration ?}


% ─────────────────────────────────────────────────────────────────────────────
\section*{Structure du Rapport}
% ─────────────────────────────────────────────────────────────────────────────

Ce rapport est organisé en quatre chapitres complémentaires, articulés de manière
progressive depuis les fondements conceptuels jusqu'à la validation expérimentale.

\textbf{Le Chapitre 1} pose les fondements du paradigme DevSecOps en retraçant
son émergence depuis le mouvement DevOps et en détaillant ses principes fondateurs,
notamment le \textit{Shift Left Security}. Il dresse ensuite un panorama structuré
des quatre grandes familles de tests de sécurité :SAST, DAST, IAST et RASP 
en mettant en évidence leurs complémentarités et leurs limites respectives.
Une attention particulière est portée au DAST dans les architectures microservices,
dont les défis spécifiques ; explosion de la surface d'attaque, dynamicité de
l'infrastructure, inaccessibilité des flux internes  constituent la problématique
centrale qui motive l'ensemble du projet.

\textbf{Le Chapitre 2} introduit le concept de Service Mesh comme réponse
architecturale aux limites identifiées au chapitre précédent. Il présente
l'architecture fondamentale du Service Mesh, structurée autour du Data Plane
et du Control Plane, ainsi que les composants clés d'Istio : \texttt{istiod},
les proxies Envoy, \texttt{istio-ingressgateway}, Kiali et la stack d'observabilité
Prometheus/Grafana. Un panorama comparatif des principales solutions du marché
(Istio, Linkerd, Consul, Cilium) justifie le choix d'Istio pour ce projet,
avant d'analyser les apports concrets du Service Mesh dans les pratiques DevSecOps.

\textbf{Le Chapitre 3} applique la méthodologie STRIDE au prototype de système
de paiement, afin de modéliser les menaces pesant sur une architecture microservices.
Pour chacune des six catégories ; Spoofing, Tampering, Repudiation, Information
Disclosure, Denial of Service et Elevation of Privilege ; les scénarios d'attaque
sont analysés en comparant le comportement du système avec et sans Service Mesh.
Le chapitre se conclut par une synthèse de douze scénarios d'attaque documentés,
mettant en correspondance chaque menace avec le mécanisme Istio responsable de
sa neutralisation.

\textbf{Le Chapitre 4} présente la validation expérimentale de l'ensemble des
analyses théoriques. Il décrit l'architecture du prototype déployé sur Minikube
avec Istio, détaille l'implémentation des six couches de sécurité  mTLS, JWT,
RBAC, SPIFFE, NetworkPolicy et TLS Gateway et documente les résultats obtenus
pour chacun des douze scénarios d'attaque. Les tableaux de bord Kiali et Grafana
sont exploités pour démontrer l'observabilité en temps réel de la posture de
sécurité, et une synthèse quantifiée mesure les apports concrets du Service Mesh
par rapport à un déploiement Kubernetes sans politiques Istio.

